\documentclass[11pt,letterpaper]{article}
\usepackage[margin=0.85in]{geometry}
\usepackage{iftex}
\ifPDFTeX
  \usepackage[T1]{fontenc}
  \usepackage[utf8]{inputenc}
  \usepackage{lmodern}
\else
  \usepackage{fontspec}
\fi
\usepackage{microtype}
\usepackage{xcolor}
\usepackage[hidelinks,unicode]{hyperref}
\usepackage{enumitem}
\usepackage{setspace}
\usepackage{fancyhdr}
\setstretch{1.08}
\setlength{\parindent}{0pt}
\setlength{\parskip}{0.72em}
\definecolor{Ink}{HTML}{172026}
\definecolor{Accent}{HTML}{B94B2C}
\pagestyle{fancy}
\fancyhf{}
\fancyhead[L]{\small\color{Accent}Statement of Purpose}
\fancyhead[R]{\small [[ school.name | tex ]]}
\fancyfoot[C]{\small\thepage}
\renewcommand{\headrulewidth}{0.4pt}
\hypersetup{pdftitle={Statement of Purpose — [[ school.name | tex ]]},pdfauthor={[[ applicant.name | tex ]]}}

\begin{document}
\color{Ink}
\begin{center}
  {\Large\bfseries Statement of Purpose\par}
  \vspace{0.35em}
  {\color{Accent}\small [[ school.program_name | tex ]] · [[ school.department | tex ]]\par}
  \vspace{0.2em}
  {\small [[ applicant.name | tex ]] · [[ applicant.target_term | tex ]] · Direct Ph.D.\par}
\end{center}
\vspace{0.4em}

I want to develop humanoid robots that can move, manipulate, and adapt as integrated physical
systems rather than as collections of isolated perception and control modules. My immediate
research goal is to learn generalizable whole-body loco-manipulation policies that remain
reliable under contact, sensory noise, and changes in embodiment. This goal grows directly from
my experience taking reinforcement-learning controllers through the complete sim-to-real loop on
Unitree G1 and H2 humanoids with Inspire dexterous hands. I have worked not only with policy
training, but also with the onboard perception, real-time software, deployment diagnostics, and
physical iteration required to make a learned controller function outside a simulator. For my
doctoral research, I hope to join a group where new learning methods and demanding real-robot
experiments can inform one another from the outset.

My undergraduate major, Big Data Management and Applications at Dalian University of Technology,
gives me a complementary rather than disconnected foundation for this agenda. Training in
mathematical modeling, probability, statistics, optimization, and large-scale data systems taught
me to formulate decisions in terms of objectives, uncertainty, constraints, and evidence. Concepts
such as Markov decision processes and statistical decision theory connect naturally to reinforcement
learning, while data-pipeline engineering connects to the less visible but decisive work of
collecting trajectories, managing simulation experiments, aligning multimodal observations, and
evaluating policy behavior at scale. I therefore see my interdisciplinary background as a route
into behavioral foundation models: large-scale simulation and robot data require both learning
algorithms and disciplined systems for producing, governing, and interpreting data. The central
lesson I carry from this preparation is the ability to convert mathematical and computational
ideas into functioning robotic experiments.

In my humanoid loco-manipulation work, I used Isaac Lab, mjlab, and MuJoCo with PyTorch to
construct training and evaluation workflows for coordinated whole-body behavior. Motion
retargeting provided a way to transform human or reference motion into embodiment-aware targets,
while reinforcement learning enabled the controller to recover behavior that was dynamically
feasible under robot constraints. ROS 2 connected policy inference with perception and control
components, and ONNX Runtime provided an explicit deployment path outside the training framework.
This combination changed how I evaluate a policy: visually plausible simulation rollouts are not
sufficient unless observation definitions, action interfaces, contact behavior, timing, and failure
recovery remain consistent when the same policy is transferred to hardware.

Working with Unitree G1/H2 platforms and Inspire hands also taught me that sim-to-real transfer is
an experimental discipline rather than a single domain-randomization setting. On Jetson Orin NX,
onboard perception and control must operate within concrete compute, latency, and integration
constraints. A discrepancy that appears minor in a training log can become a timing mismatch,
an unstable contact transition, or a systematic action error on the physical robot. My response
has been to trace the complete observation-to-action path, isolate where simulation and hardware
assumptions diverge, and iterate across policy representation, interfaces, and deployment tests.
I do not view this hardware work as separate from research novelty. It exposes failure modes that
can motivate better state representations, training distributions, safety mechanisms, and
evaluation protocols. It also prepares me to contribute to a laboratory from the first experiment:
I understand the difference between a policy that merely runs and one whose behavior can be
systematically explained.

[% if school.track_type == "CS" %]
Within a Computer Science track, I want to emphasize large-scale reinforcement learning,
representation learning, and visuomotor policies for general-purpose embodied behavior. A central
question is how to learn state and action representations that preserve task-relevant visual and
contact information while transferring across camera viewpoints, objects, scenes, and robot
embodiments. I am interested in pretraining behavioral representations from heterogeneous motion,
simulation, and robot trajectories, then adapting them efficiently to new loco-manipulation tasks.
My data-systems preparation is directly relevant to building reproducible training pipelines and
diagnosing distribution shift, while my deployment experience gives me a concrete standard for
whether a supposedly general representation captures something physically useful. I also want to
study evaluation beyond aggregate success: policies should be compared across visual variation,
contact regimes, latency conditions, and recovery behavior so that progress toward robot foundation
models is measurable rather than rhetorical.
[% elif school.track_type == "ME" %]
Within a Mechanical Engineering track, I want to place whole-body dynamics, contact-aware control,
and experimentally validated model learning at the center of my research. Large-scale learning can
produce expressive behaviors, but humanoid loco-manipulation still depends on force interaction,
stability, actuator limits, and the mechanics of hands and bodies. I am interested in combining
optimization or model-based structure with learned residual dynamics and adaptive policies, so that
data improves performance without discarding physical interpretation. My experience with humanoid
and dexterous-hand deployment would allow me to test these ideas through controlled comparisons of
contact transitions, disturbances, and changes in load. I would approach sim-to-real transfer as a
way to identify and reduce modeling error, connecting simulation fidelity, system identification,
control design, and repeatable mechanical experiments.
[% else %]
Within a robotics-centered track, I want to integrate whole-body dynamics, contact-rich control,
perception, and hardware deployment with scalable learning. Humanoid loco-manipulation offers an
ideal setting for this integration: locomotion cannot be separated from manipulation when both
compete for balance, contact, and limited onboard computation. I hope to investigate policies that
coordinate feet, arms, hands, and gaze while adapting to uncertainty in object properties and the
surrounding environment. I am particularly interested in combining model-based constraints with
learned representations, using motion retargeting to structure behavior without restricting the
controller to a single embodiment. My end-to-end deployment background would help me design
experiments in which algorithms, sensing, software timing, and robot mechanics are evaluated as a
single coupled system.
[% endif %]

The research direction I hope to pursue in my Ph.D. is a hardware-grounded behavioral foundation
model for humanoid loco-manipulation. Such a model should not simply absorb larger trajectory
datasets; it should support purposeful adaptation, expose uncertainty, and preserve the physical
structure needed for safe execution. I imagine a training framework in which motion retargeting
and simulation provide diverse structured experience, real-robot trials identify systematic gaps,
and targeted fine-tuning converts those gaps into new learning signals. Vision and proprioception
should be learned jointly with action so that a policy can anticipate contact and recover when its
expectations are violated. Equally important, I want to build evaluation suites that report
generalization, robustness, recovery, and deployment cost alongside task success. This agenda
connects my interests in reinforcement learning, sim-to-real transfer, visual perception, and
robot foundation models while giving each component a testable role.

[[ school.why_school | tex ]] [[ school.why_track | tex ]]

[% for pi in target_pis[:2] %]
I am particularly interested in Professor [[ pi.name | tex ]]'s research in
[[ pi.recent_work_keywords[0] | tex ]] and [[ pi.recent_work_keywords[1] | tex ]]. I would bring
a concrete experimental perspective to this connection: my closed-loop humanoid and dexterous-hand
experience would allow me to test new representations and control methods under real contact,
sensing, and computation constraints. The question I hope to pursue is how progress in
[[ pi.recent_work_keywords[0] | tex ]] can support adaptable behavior whose generalization is
demonstrated through repeatable physical interaction, not only through simulation benchmarks.

[% endfor %]
My long-term objective is to [[ profile.sop.career_goal | tex ]] I want my doctoral training to
make me equally comfortable deriving a learning objective, examining a contact model, building a
data pipeline, and diagnosing a failed hardware trial. [[ school.name | tex ]]'s
[[ school.department | tex ]] program is compelling to me because it would challenge all of these
capabilities rather than force me to choose between algorithms and physical systems. With my
interdisciplinary preparation and hands-on humanoid experience, I am ready to contribute sustained
experimental work, collaborative engineering, and a research agenda aimed at capable and reliable
robots in the real world.

\end{document}
